% Authority: German Wikiversity pageid 131353, revid 1111802, 2026-08-15T16:41:12Z.
% Exact UTF-8 witness SHA-256: b07aac01b6f803481fab9b5331e19480fb9bf0058597159b6152c65c7bc955ce.

\zwischenueberschrift{Solusi untuk Soal 1}

Berlaku

\mavergleichskettedisp
{\vergleichskette
{ v
}
{ \in }{ T_PZ
}
{ = }{ \ker \left((D\varphi)_P\right)
}
{ } {
}
{ } {
}
}
{}{}{.}
Menurut
\definitionsverweis {teorema fungsi implisit}{}{,}
terdapat himpunan terbuka

\mathbed {P \in W} {}
 {W \subseteq G} {}
 {} {} {} {,}
himpunan terbuka

\mavergleichskette
{\vergleichskette
{ V
}
{ \subseteq }{ \R^{n-m}
}
{ } {
}
{ } {
}
{ } {
}
}
{}{}{,}
dan pemetaan terdiferensialkan secara kontinu
\maabbdisp {\psi} { V } { W } {}
sedemikian sehingga

\mavergleichskette
{\vergleichskette
{ \psi(V)
}
{ \subseteq }{ Z \cap W
}
{ } {
}
{ } {
}
{ } {
}
}
{}{}{}
dan $\psi$ menginduksi
\definitionsverweis {bijeksi}{}{}
\maabbdisp {\psi} { V } { Z \cap W } {}.
Misalkan

\mavergleichskette
{\vergleichskette
{ Q
}
{ \in }{ V
}
{ } {
}
{ } {
}
{ } {
}
}
{}{}{}
adalah titik yang memenuhi

\mavergleichskettedisp
{\vergleichskette
{ \psi(Q)
}
{ =} { P
}
{ } {
}
{ } {
}
{ } {
}
}
{}{}{.}
Pemetaan $\psi$ di setiap titik

\mavergleichskette
{\vergleichskette
{ Q
}
{ \in }{ V
}
{ } {
}
{ } {
}
{ } {
}
}
{}{}{}
bersifat
\definitionsverweis {reguler}{}{}
dan
\definitionsverweis {diferensial total}{}{}
dari $\psi$ memenuhi

\mavergleichskettedisp
{\vergleichskette
{ (D\varphi)_P \circ (D\psi)_Q
}
{ =} { 0
}
{ } {
}
{ } {
}
{ } {
}
}
{}{}{.}
Dengan demikian,

\mavergleichskettedisp
{\vergleichskette
{ \operatorname{im} \left((D\psi)_Q\right)
}
{ =} { \ker \left((D\varphi)_P\right)
}
{ =} { T_PZ
}
{ } {
}
{ } {
}
}
{}{}{.}
Karena $\psi$ reguler di $Q$, pemetaan
\maabbdisp {(D\psi)_Q} { \R^{n-m} } { \R^n } {}
bersifat injektif, sedangkan pemetaan
\maabbdisp {(D\psi)_Q} { \R^{n-m} } { T_PZ } {}
bersifat bijektif. Misalkan

\mavergleichskette
{\vergleichskette
{ u
}
{ \in }{ \R^{n-m}
}
{ } {
}
{ } {
}
{ } {
}
}
{}{}{}
adalah prapeta dari $v$, dan misalkan
\maabbeledisp {\theta} { ]-\epsilon,\epsilon[ } { \R^{n-m} } {t} {Q+tu} {,}
dengan $\epsilon$ dipilih cukup kecil agar seluruh citra terletak di dalam $V$. Maka
\maabbdisp {\gamma=\psi\circ\theta} { ]-\epsilon,\epsilon[ } { \R^n } {}
mempunyai sifat

\mavergleichskettedisp
{\vergleichskette
{ \gamma(0)
}
{ =} { \psi(\theta(0))
}
{ =} { \psi(Q)
}
{ =} { P
}
{ } {
}
}
{}{}{}
dan

\mavergleichskettedisp
{\vergleichskette
{ \gamma'(0)
}
{ =} { (D\psi)_Q(\theta'(0))
}
{ =} { (D\psi)_Q(u)
}
{ =} { v
}
{ } {
}
}
{}{}{.}
